% Reviewer 2 standalone response; common front matter and R2 content from response-to-reviewers.tex.
\documentclass[11pt]{article}

\usepackage[a4paper,margin=24mm]{geometry}
\usepackage{fontspec}
\usepackage{xcolor}
\usepackage[hidelinks]{hyperref}

\setmainfont[
  Path=../latex-kiee-review-2022/fonts/,
  Extension=.otf,
  UprightFont=*-Regular,
  BoldFont=*-Bold
]{NotoSerifCJKkr}

\definecolor{reviewercolor}{RGB}{38,82,126}
\definecolor{responsecolor}{RGB}{36,113,78}
\definecolor{revisioncolor}{RGB}{145,82,24}

\setlength{\parindent}{0pt}
\setlength{\parskip}{0.55em}
\setlength{\fboxsep}{7pt}
\setlength{\emergencystretch}{3em}
\renewcommand{\familydefault}{\rmdefault}

\newcommand{\fieldheading}[2]{%
  \par\medskip\noindent
  {\color{#1}\rule{3pt}{1.25em}}\hspace{0.55em}\textbf{#2}\par\smallskip}

\newenvironment{reviewitem}[1]
  {\subsection*{#1}}
  {\par\medskip\hrule\bigskip}
\newenvironment{reviewercomment}
  {\fieldheading{reviewercolor}{1. 리뷰어 코멘트}%
   \begin{quote}}
  {\end{quote}}
\newenvironment{authorresponse}
  {\fieldheading{responsecolor}{2. 저자 답변}}
  {\par}
\newenvironment{manuscriptrevision}[1]
  {\fieldheading{revisioncolor}{3. 논문 반영 결과}%
   \textbf{수정 위치:} #1\par\smallskip}
  {\par}

\newcommand{\manuscripttitle}{비전-언어 자율주행 모델의 실행 가드 기반 궤적 후보 선택: 안전 제약 인과 연쇄}
\newcommand{\manuscriptid}{2026-D-D04-0036}
\newcommand{\revisionround}{1차 수정}
\newcommand{\responsedate}{2026년 9월 4일}

\begin{document}

\begin{center}
  {\LARGE\bfseries 심사 의견 답변서\par}
  \vspace{1em}
  {\large\bfseries \manuscripttitle\par}
\end{center}

\vspace{0.8em}
\begin{tabular}{@{}ll@{}}
논문 접수번호: & \manuscriptid \\
수정 차수: & \revisionround \\
답변일: & \responsedate
\end{tabular}

\section*{답변에 앞서}

편집위원장님과 리뷰어 여러분께,

원고를 세심하게 검토하고 건설적인 의견을 주셔서 감사합니다. 저자들은 모든 의견을
검토하여 원고를 수정했습니다. 아래에는 각 리뷰어 코멘트, 이에 대한 저자 답변,
그리고 원고에 실제로 반영한 내용과 수정 위치를 순서대로 제시합니다. 쪽 번호는
개정 원고를 기준으로 합니다.

이번 개정에서는 후보 정확도와 계약 쌍 정확도의 관계를 명확히 하고, 제약의 의미와
입력 위치의 효과를 분리하여 설명했습니다. 또한 고난도 사례와 결정적 단위 정규화
분석을 추가하고 실행 차단장치의 개입 비용을 보고했습니다. 기존 주행 비전-언어
모델 및 실행시간 검증 연구와의 차이, 국문 요약과 영문 초록의 구체성, 원고의 논리적
흐름, 학습된 가드 준수와 실제 실행 안전의 구분도 보강했습니다. 모델 외부 검증과
평가 지표의 상호 보완적 역할도 명시했습니다. 리뷰어 1에 따른 수정은 파란색,
리뷰어 2에 따른 수정은 보라색으로 표시했습니다.

\section*{리뷰어 2}

\begin{reviewitem}{R2-C1}
  \begin{reviewercomment}
    기존 비전-언어 자율주행 모델과 본 연구의 차별성 및 기존 연구의 문제점에 대한
    설명이 부족하며 연구동기가 불분명함. 즉 본연구의 필요성 및 기존 연구의
    문제점에 대한 설명이 부족하여 연구 동기가 불분명함.
  \end{reviewercomment}

  \begin{authorresponse}
    지적에 감사드립니다. 개정 원고에서는 기존 연구를 두 흐름으로 나누어 연구
    공백을 명확히 했습니다. Driving VLM/CoC는 장면의 원인, 행동 이유 및 목표를
    설명하지만 동일 목표 후보의 허용 경계만 바꾸었을 때 선택이 전환되는지를
    직접 평가하지 않습니다. 반면 형식 검증기와 runtime monitor는 주어진 수치
    명세를 정확히 검사하지만 자연어 가드를 VLM의 학습 입력으로 사용하고 그
    의미 반영 여부를 측정하지 않습니다. 본 연구는 기존 계획기나 검증기를
    대체한다고 주장하지 않습니다. CoC에 자연어 실행 가드를 결합하고, 장면ㆍ목표ㆍ
    후보를 고정한 paired-contract와 모델 외부 수치 검증을 연결하는 중간
    인터페이스를 제안한다는 차별점을 명시했습니다.
  \end{authorresponse}

  \begin{manuscriptrevision}{제1절(1--2쪽), 제2.4절과 표 2(3--4쪽)}
    서론에 두 연구 흐름 사이의 공백과 필요한 인터페이스를 설명하는 동기 문단을
    추가했습니다. 관련 연구에는 driving VLM/CoC, formal monitor/shield 및 본
    연구의 주요 역할, 실행 조건, 진단 평가와 증거 범위를 비교하는 표를
    추가했습니다. 해당 개정은 보라색으로 표시했습니다.
  \end{manuscriptrevision}
\end{reviewitem}

\begin{reviewitem}{R2-C2}
  \begin{reviewercomment}
    전체적인 구성에 있어 논리적 연결성이 미비함. 전체적으로 논리적인 흐름을
    정리할 필요가 있음.
  \end{reviewercomment}

  \begin{authorresponse}
    의견에 따라 원고의 논리적 연결을 ``연구 공백--문제 정의--입력 및 평가 방법--
    단계별 실험--증거 범위''의 순서로 명시했습니다. 서론 말미에 절별 roadmap을
    추가하고, 문제 정의에서는 행동 목표와 가드 만족을 분리한 이유를 먼저
    설명했습니다. 방법 절은 입력, 대조군과 paired-contract, 모델 외부 판정의 세
    질문으로 구성했으며, 실험 절은 사전 진단에서 정적ㆍ시간ㆍ다중 주행 및 별도
    폐루프 분석으로 범위를 넓히는 순서를 명시했습니다. 각 실험은 확인 목적,
    비교 방법, 결과 및 의미의 동일한 구조로 읽히도록 점검했습니다.
  \end{authorresponse}

  \begin{manuscriptrevision}{제1절(2쪽), 제3절(3쪽), 제4절(4쪽), 제5절(7쪽)}
    절별 roadmap과 문제 정의ㆍ방법ㆍ실험 절의 전환 문단을 추가했습니다. 이
    문단들은 앞 절의 질문이 다음 절의 수식, 평가 절차 및 실험 순서로 어떻게
    이어지는지 명시합니다. 해당 개정은 보라색으로 표시했습니다.
  \end{manuscriptrevision}
\end{reviewitem}

\begin{reviewitem}{R2-C3}
  \begin{reviewercomment}
    국영문 요약에 제안한 방법에 대한 구체적인 내용을 기술하시기 바랍니다.
  \end{reviewercomment}

  \begin{authorresponse}
    국문 요약을 추가하고 영문 초록을 전면 보완했습니다. 두 요약 모두 (1) 기존
    CoC에서 빠질 수 있는 실행 허용 범위, (2) 일곱 요소의 자연어 execution guard,
    (3) 장면ㆍCoCㆍ가드ㆍ수치 후보를 입력받는 선택 절차, (4) 모델 외부 결정적
    검증, (5) guard만 바꾸는 paired-contract 평가를 구체적으로 기술합니다.
    핵심 결과로 표준 시간 평가의 위반률 0.319에서 0.000으로의 감소와 준수 목표
    완료율/paired accuracy 1.000, 미관측 시간값의 후보 정확도 0.635, 위반률
    0.351 및 paired accuracy 0.306을 함께 제시해 효과와 한계를 균형 있게
    요약했습니다.
  \end{authorresponse}

  \begin{manuscriptrevision}{국문 요약 및 영문 초록(1쪽)}
    국문 요약을 새로 추가하고 영문 초록을 동일한 구성으로 개정했습니다. 제안
    방법의 입력, 후보 선택, 모델 외부 검증과 쌍대 계약 평가를 구체화했습니다.
    핵심 수치와 ``실제 차량 안전 보장이 아님''이라는 주장 범위도 일관되게
    기술했습니다. 해당 개정은 보라색으로 표시했습니다.
  \end{manuscriptrevision}
\end{reviewitem}

\begin{reviewitem}{R2-C4}
  \begin{reviewercomment}
    학습된 가드 준수와 실제 실행안전의 구분을 한 이유가 무엇인지 구체적인 기술이
    필요합니다.
  \end{reviewercomment}

  \begin{authorresponse}
    두 개념을 구분한 이유는 평가 단위와 실패 원인이 다르기 때문입니다. 본 논문의
    learned guard compliance는 유한한 합성 후보와 연구자가 제공한 부분 계약에서
    목표와 가드를 만족하는 후보를 선택한 비율입니다. 실제 execution safety는
    센서 오인식, 계약에서 빠진 위험, 연속 궤적과 차량 동역학의 오차, 다른 도로
    사용자의 반응 및 지연된 상태 갱신까지 포함하는 폐루프 성질입니다. 따라서
    후보 준수율이 1.000이어도 가드가 불완전하거나 관측이 틀리면 실제 차량은
    위험할 수 있습니다. 이 구분은 합성 평가 점수를 안전 보장으로 확대하는 범주
    오류를 방지하고, 최신 상태를 사용하는 외부 검증기와 실행 차단장치가 별도로
    필요한 이유를 설명합니다.
  \end{authorresponse}

  \begin{manuscriptrevision}{제4.6절(7쪽), 제6.4절(13--14쪽)}
    학습된 가드 준수와 실제 실행 안전의 평가 단위 및 포함 실패 원인을 명시하고,
    두 개념을 구분해야 하는 이유를 추가했습니다. 폐루프 논의에서는 계약 상태의
    가역성과 최신성, 실행 차단의 역할을 이 구분과 연결했습니다. 해당 개정은
    보라색으로 표시했습니다.
  \end{manuscriptrevision}
\end{reviewitem}

\begin{reviewitem}{R2-C5}
  \begin{reviewercomment}
    연구 내용에 정확한 정보 전달을 위해서 문장 구성에 대한 점검도 필요합니다.
  \end{reviewercomment}

  \begin{authorresponse}
    원고 전체에서 용어와 문장 주체를 점검했습니다. ``기본 CoC'', ``제약 통합
    CoC'', ``실행 가드'', ``모델 외부 결정적 검증기'', ``가드 준수'' 및 ``실제
    실행 안전''을 일관되게 사용했습니다. 특히 연구 공백, 방법 입력ㆍ출력, 검증
    주체 및 결과의 해석 범위를 한 문장에 혼합하지 않도록 나누었습니다. 또한 각
    실험을 확인 목적, 비교 방법, 관찰 결과 및 의미로 분리하고, 서로 다른 모델의
    결과가 직접 비교되지 않는다는 주어와 범위를 명시했습니다.
  \end{authorresponse}

  \begin{manuscriptrevision}{제1--5절(1--7쪽), 제6.1절(13쪽)}
    대표적으로 서론의 연구 공백을 두 선행 연구 흐름과 본 연구의 역할로 분리하고,
    방법 절을 ``입력--대조 평가--외부 판정''의 세 질문으로 다시 설명했습니다.
    실험 절에는 증거의 순서를 안내하는 문단을 추가했으며, 논의에서는 관찰된
    수치와 허용 가능한 해석을 별도 문장으로 구분했습니다. 주요 보완 문장은
    보라색으로 표시했습니다.
  \end{manuscriptrevision}
\end{reviewitem}

\begin{reviewitem}{R2-C6}
  \begin{reviewercomment}
    제안하는 방법에 대해 또한 어떠한 장점이 있는지를 기술한다면 보다 완성도가
    높을 것 입니다.
  \end{reviewercomment}

  \begin{authorresponse}
    제안 방법의 장점을 네 가지로 명시했습니다. 첫째, 기존 CoC의 설명 구조를
    유지하면서 허용ㆍ금지ㆍ해제 조건을 자연어로 추가합니다. 둘째, paired-contract
    평가로 장면 암기와 guard에 따른 선택 전환을 구분합니다. 셋째, violation뿐
    아니라 goal completion, deadlock 및 excess conservatism을 함께 측정하여 항상
    정지하는 해법을 성공으로 보지 않습니다. 넷째, 모델의 후보 선택과 결정적 수치
    판정을 분리해 외부 verifier나 shield로 연결할 수 있습니다. 이러한 장점은
    합성 후보 선택의 해석 가능성과 진단 가능성에 관한 것이며 실제 차량의 안전
    보장으로 확대하지 않았습니다.
  \end{authorresponse}

  \begin{manuscriptrevision}{제1절(2쪽), 제6.1절(13쪽)}
    서론에 네 가지 장점을 명시하고, 논의에서 이 장점이 범용 주행 모델의 우월성이
    아니라 허용조건의 명시성, 반사실적 진단 및 독립 수치 판정에 있음을 기존
    CoC-only 조건과 대비했습니다. 해당 개정은 보라색으로 표시했습니다.
  \end{manuscriptrevision}
\end{reviewitem}

\begin{reviewitem}{R2-C7}
  \begin{reviewercomment}
    또한, 모델 외부 궤적 검증과 평가지표를 이용해서 제안하는 방법에 대한 성능이
    우수하다는 부분을 부각시킬 필요가 있습니다.
  \end{reviewercomment}

  \begin{authorresponse}
    외부 trajectory verifier는 모델의 자기 설명을 신뢰하지 않고 후보에 저장된
    정지 위치, 속도, 감속도, 가가속도, 간격 및 진입 시각을 guard 경계와 직접
    비교합니다. 이를 통해 candidate accuracy, violation, guard-compliant goal
    completion, deadlock, excess conservatism 및 paired-contract accuracy를 상호
    보완적으로 계산합니다. 표준 정적ㆍ시간ㆍ정지/차선 변경 평가에서 제약 통합
    CoC의 violation rate는 모두 0.000, paired-contract accuracy는 모두 1.000이었고
    목표 완료 관련 지표도 1.000이었습니다. 따라서 낮은 위반률이 동일 답 반복이나
    목표 포기의 결과가 아님을 강조했습니다. 동시에 미관측 시간값에서 위반률
    0.351, paired accuracy 0.306으로 악화된 결과를 함께 제시하여 우수성의 범위를
    표준 합성 평가로 한정했습니다.
  \end{authorresponse}

  \begin{manuscriptrevision}{제4.5절(6--7쪽), 제6.1절(13쪽), 제7절(15쪽)}
    방법 절에 외부 검증과 각 지표의 상호 보완적 진단 역할을 추가했습니다. 논의와
    결론에서는 표준 평가의 violation/goal completion/paired accuracy를 함께
    해석하여 제안 방법의 성능을 부각하고, 미관측 입력 결과를 근거로 적용 범위를
    제한했습니다. 해당 개정은 보라색으로 표시했습니다.
  \end{manuscriptrevision}
\end{reviewitem}

\section*{맺음말}

원고의 명확성과 엄밀성을 높일 수 있도록 세심한 의견을 주신 리뷰어 여러분께 다시
한번 감사드립니다. 개정 원고와 본 답변서가 제기된 모든 우려에 충실히 답하였기를
바랍니다.

감사합니다.\\
저자 일동

\end{document}
